% !PRINTED 259
\weberpage{259}

\begin{center}
\S.\,67.

\gesperrt{Congruence Fields of the Second Degree.}
\end{center}

In every congruence field $\mathfrak{E}_{n,p}$ for the modulus $p$ of degree
$n$, the congruence field of the first degree $\mathfrak{E}_{1,p}$, which
consists of the rational numbers taken with respect to the modulus $p$, is
contained as a divisor. Hence in the group of linear
substitutions $E_{n,p}$ there is also contained as a divisor a group which
consists of all substitutions of the form

\weq{1}{A = \begin{pmatrix} a, & b \\ c, & d \end{pmatrix}, \quad ab-cd=1 % [sic] printed ab-cd, not ad-bc
}
where $a, b, c, d$ are integral rational numbers taken with respect to the
modulus $p$.

This group, which appears as the group of the modular equations
in the theory of elliptic functions$^{1)}$, is the
true goal of our considerations. But for the
investigation of this group, and especially for the search for
its divisors, it offers the greatest advantage not to consider this group
independently by itself, but as a divisor of a
group $E_{2,p}$. Namely, if we denote by $N$ any fixed
quadratic non-residue of $p$, then, as already remarked
above, $t^2-N$ is irreducible with respect to the modulus $p$, and we obtain
a congruence field of degree $2$ when, by the equation

\weq{2}{\varepsilon^2 = N}
we introduce a \gesperrt{Galois} imaginary number.

For the sake of clarity it shall be stipulated for the next
considerations that the small Latin letters denote
integral rational numbers (taken with respect to the modulus $p$), which we
here also call \gesperrt{real numbers}, and the Greek letters
numbers of the field $\mathfrak{E}_{2,p}$.

This field $\mathfrak{E}_{2,p}$ has the property, very important for us,
that \gesperrt{in it all real numbers are squares}.

For if $a$ is a quadratic residue of $p$, we can
set $a=x^2$; and if $b$ is a quadratic non-residue, then
$x^2=bN$ can be satisfied, hence $b=(x\varepsilon^{-1})^2$.

\rule{4cm}{0.4pt}

$^{1)}$ Compare the author's book: ``Elliptische Functionen und
algebraische Zahlen''. Braunschweig 1891.
